\section*{جمع‌بندی اهداف طراحی و هزینه}
\phantomsection
\addcontentsline{toc}{section}{جمع‌بندی اهداف طراحی و هزینه}
\label{sec:book-objectives}
جمع‌بندی فصل ۱ کتاب چهار هدف \lr{Reliability}، \lr{Performance}، \lr{Power}
و \lr{Cost} را کنار هم قرار می‌دهد \cite[صفحات چاپی ۵۵ و ۵۶]{weste}.
\lr{Reliability} به عملکرد قابل‌قبول در شرایط و طول عمر مورد انتظار مربوط است.
ولتاژ و دمای مجاز، بخشی از این شرایط‌اند.
\lr{Performance} به زمان انجام کار و توان عملیاتی وابسته است و با فرکانس
\lr{Clock} و میزان موازی‌سازی ارتباط دارد.
توان، محدودیت حرارتی و مصرف انرژی را به انتخاب مدار و فعالیت آن پیوند می‌دهد.
هزینه نیز از تلاش طراحی، ماسک و تولید هر قطعه اثر می‌پذیرد.

\subsection*{هزینه ثابت و هزینه هر قطعه}
\lr{Nonrecurring Engineering} یا \lr{NRE} هزینه‌هایی مانند طراحی و تهیه ماسک
را دربرمی‌گیرد. برای مدل اقتصادی آموزشی، \(K\) را هزینه ثابت، \(c\) را هزینه
تولید هر قطعه و \(N>0\) را تعداد قطعه می‌نامیم.
با فرض ثابت‌بودن \(c\) و لحاظ‌شدن شرایط تولید در همین ضریب، هزینه کل و میانگین برابرند با
\begin{equation}
 C_{\mathrm{total}}=K+Nc,\qquad
 C_{\mathrm{unit}}=\frac{K}{N}+c.
 \label{eq:booklet-nre}
\end{equation}
با افزایش تعداد تولید، سهم هزینه ثابت در هر قطعه کاهش می‌یابد.
کتاب این سازوکار را در مقایسه فناوری با هزینه مهندسی بالا و فناوری بالغ‌تر
توضیح می‌دهد. هزینه واقعی به مساحت، فرآیند، \lr{Yield}، بسته‌بندی و شرایط تولید
وابسته است. اعداد نمونه این پاراگراف برای تحلیل روش فرض می‌شوند.

دو انتخاب فرضی \(A\) و \(B\) با داده‌های
\(K_A=100000\)، \(c_A=5\)، \(K_B=400000\) و \(c_B=2\)، همگی در واحد پول
سازگار، را در نظر بگیرید.
برای \(N=10000\)، میانگین‌ها \(100000/10000+5=15\) و
\(400000/10000+2=42\) واحد پول‌اند.
برای یافتن تعداد برابری، دو رابطه را برابر می‌کنیم
\begin{equation}
 \frac{100000}{N}+5=\frac{400000}{N}+2
 \quad\Longrightarrow\quad
 3=\frac{300000}{N}
 \quad\Longrightarrow\quad N=100000.
 \label{eq:booklet-cost-crossover}
\end{equation}
در این تعداد، هر دو میانگین ۶ واحد پول می‌شوند.
برای تولید کمتر، \(A\) و برای تولید بیشتر، \(B\) در همین فرض‌ها هزینه میانگین
کمتری دارد. این مثال از مفهوم هزینه ثابت در جمع‌بندی کتاب بسط داده شده است.

\subsection*{پیوند مدل‌های درس با تصمیم طراحی}
انتخاب شبکه منطقی و اندازه ترانزیستور، جریان و ظرفیت را تغییر می‌دهد.
این دو کمیت، زمان شارژ و تخلیه را تعیین می‌کنند.
انتخاب \lr{Layout}، طول سیم، مساحت پیوندها و فضای سیم‌کشی را تغییر می‌دهد.
روش ساخت و بسته‌بندی نیز بر محدودیت الکتریکی، هزینه و قابلیت اطمینان اثر دارند.
برای ارزیابی یک تصمیم، باید فرض عملکرد، معیار زمان، بودجه توان و مدل فیزیکی
مورد استفاده مشخص باشد.
فصل ۲ این جزوه، مدل جریان و ظرفیت لازم برای آغاز این محاسبه را فراهم می‌کند.
\clearpage
